% RQ2 executability evaluation snippet.
% Required package: \usepackage{booktabs}

\subsection{Executability of Generated Robot Framework Test Cases}

RQ2 evaluates whether the generated Robot Framework test cases can be executed
successfully against the running SUT. Two datasets were used: 17 Gherkin
scenarios from \texttt{requirements/gherkin} and 30 natural-language
requirements from \texttt{requirements/userstories}. The natural-language
dataset contains repeated intents expressed with different wording, because the
SUT does not implement enough distinct features to provide one unique feature
for every requirement. Therefore, paraphrased requirements are kept in the
dataset to test whether the approach remains robust under wording variation.

Each generated test case was executed against the SUT at
\texttt{http://localhost:4242}. Failures were classified into three main
categories. Mapping errors indicate that the generated test selected a wrong or
unsupported keyword, locator, route, or UI element. Logic errors indicate that
the generated steps were syntactically valid but ordered incorrectly for the
required SUT state, for example accessing a protected page without first
logging in. Actual SUT defects indicate that the generated test was executable
but failed because the SUT did not provide the expected behavior or element.

\begin{table}[htbp]
\centering
\caption{RQ2 executability results for Gherkin requirements.}
\label{tab:rq2_gherkin_executability}
\begin{tabular}{llrrrrrr}
\toprule
Pipeline & Type & Total & Passed & Success (\%) & Mapping & Logic & SUT Defects \\
\midrule
Baseline & Gherkin Scenario & 17 & 0 & 0.00 & 17 & 0 & 0 \\
Proposed & Gherkin Scenario & 17 & 17 & 100.00 & 0 & 0 & 0 \\
\bottomrule
\end{tabular}
\end{table}

\begin{table}[htbp]
\centering
\caption{RQ2 executability results for natural-language requirements.}
\label{tab:rq2_natural_executability}
\begin{tabular}{llrrrrrr}
\toprule
Pipeline & Type & Total & Passed & Success (\%) & Mapping & Logic & SUT Defects \\
\midrule
Baseline & Bug Report & 2 & 0 & 0.00 & 0 & 2 & 0 \\
Baseline & Functional Requirement & 11 & 6 & 54.55 & 0 & 5 & 0 \\
Baseline & User Story & 17 & 3 & 17.65 & 0 & 14 & 0 \\
Proposed & Bug Report & 2 & 1 & 50.00 & 0 & 0 & 1 \\
Proposed & Functional Requirement & 11 & 7 & 63.64 & 0 & 0 & 3 \\
Proposed & User Story & 17 & 16 & 94.12 & 0 & 0 & 1 \\
\bottomrule
\end{tabular}
\end{table}

\begin{table}[htbp]
\centering
\caption{Numbered proposed-approach failures in the natural-language dataset.}
\label{tab:rq2_proposed_natural_failures}
\begin{tabular}{llll}
\toprule
Test Case & Requirement Type & Classification & Failure Cause \\
\midrule
N013 & Functional Requirement & Actual SUT Defect & Remember Me expected active, actual state inactive \\
N018 & Functional Requirement & Actual SUT Defect & Expected lesson/course action button not available \\
N019 & User Story & Actual SUT Defect & Expected lesson/course action button not available \\
N023 & Functional Requirement & Actual SUT Defect & Expected exercise submit flow not available \\
N026 & Bug Report & Actual SUT Defect & Expected wrong-answer heart behavior not executable as required \\
N028 & Functional Requirement & Non-UI & Password hashing requires backend/database assertion \\
\bottomrule
\end{tabular}
\end{table}

The proposed approach executed all 17 generated Gherkin scenarios successfully,
achieving a 100.00\% execution success rate. The zero-shot baseline failed all
Gherkin scenarios because it generated generic or unsupported Robot/Selenium
style keywords instead of project-valid Robot Framework keywords, resulting in
17 mapping failures.

For natural-language requirements, the proposed approach achieved an execution
success rate of 80.00\% overall. It performed strongest on user stories, where
16 of 17 generated tests passed (94.12\%). Functional requirements were more
difficult because several requirements refer to behavior that is missing or only
partially implemented in the SUT, such as lesson exercise actions and checkbox
state expectations. The zero-shot baseline achieved only 30.00\% overall
success on the same natural-language dataset. Most baseline failures were logic
errors caused by generic selectors, ambiguous text matching, or attempts to
access protected pages without applying the required authenticated workflow.

These results show that the proposed approach substantially improves
executability compared with the baseline. The main remaining failures are not
caused by Robot Framework syntax errors or keyword hallucination; instead, they
are concentrated in SUT limitations, missing UI behavior, or requirements that
cannot be validated through the browser UI alone.
